\subsection{Interpretation of the product measure as an integral of measures}

Throughout this section, T and \(T'\) denote two locally compact spaces, \(\mu\) a positive measure on T, \(\mu'\) a positive measure on \(T'\) , and \(\nu = \mu \otimes \mu'\) the product measure on \(X = T \times T'\) (Ch. III, §4, No.~1).

For every \(t \in T\) , the mapping \(t' \mapsto (t, t')\) of \(T'\) into X is continuous and proper. Let \(\lambda_t'\) be the image of \(\mu'\) under this mapping; \(\lambda_t'\) is a positive measure on X, and if \(f \in \mathcal{K}(X)\) then, denoting by \(f_t\) the partial mapping \(t' \mapsto f(t, t')\) , we have

\[
\int f d \lambda_ {t} ^ {\prime} = \int f _ {t} d \mu^ {\prime},\tag{1}
\]

which is also expressed by the relation \(\lambda_t' = \varepsilon_t \otimes \mu'\) .

Moreover, the mapping \(t \mapsto \lambda_t'(f)\) is continuous, with compact support (Ch. III, §4, No.~1, Lemma 2), therefore the mapping \(t \mapsto \lambda_t'\) of T into \(\mathcal{M}(\mathbf{X})\) is vaguely continuous (and, a fortiori, vaguely \(\mu\) -measurable); consequently, the family of measures \(t \mapsto \lambda_t'\) is \(\mu\) -adequate (§3, No.~1, Prop. 2a)). The integral of \(f\) with respect to the measure \(\int \lambda_t' d\mu(t)\) is by definition

\[
\int \langle f, \lambda_ {t} ^ {\prime} \rangle d \mu (t) = \int d \mu (t) \int f _ {t} (t ^ {\prime}) d \mu^ {\prime} (t ^ {\prime}) = \int f (t, t ^ {\prime}) d \nu (t, t ^ {\prime})
\]

(Ch. III, §4, No.~1, Th. 2); thus \(\nu = \int \lambda_t' d\mu(t)\) .

Similarly, for every element \(t^{\prime} \in T^{\prime}\) , let \(\lambda_{t^{\prime}}\) be the image of \(\mu\) under the mapping \(t \mapsto (t, t^{\prime})\) of T into X. The mapping \(t^{\prime} \mapsto \lambda_{t^{\prime}}\) is \(\mu^{\prime}\) -adequate and vaguely continuous, and \(\nu = \int \lambda_{t^{\prime}} d\mu^{\prime}(t^{\prime})\) . We shall need the following lemmas:

Lemma 1. --- For every numerical function \(f \geqslant 0\) defined on \(X\) ,

\[
\int^ {*} f _ {t} d \mu^ {\prime} = \int^ {*} f d \lambda_ {t} ^ {\prime}.\tag{2}
\]

Since \(t' \mapsto (t, t')\) is a continuous and proper mapping, this follows from Prop. 2 of §4, No.~2.

Lemma 2. --- Let \(f\) be a mapping of \(X\) into a topological space. For \(f\) to be \(\lambda_t'\) -measurable, it is necessary and sufficient that \(f_t\) be \(\mu'\) -measurable. This is a consequence of Prop. 3 of §6, No.~2.

Lemma 3. --- Let f be a function defined on X, with values in \(\overline{R}\) or in a Banach space. For f to be \(\lambda_{t}^{\prime}\) -integrable, it is necessary and sufficient that \(f_{t}\) be \(\mu^{\prime}\) -integrable, in which case

\[
\int \mathbf {f} _ {t} d \mu^ {\prime} = \int \mathbf {f} d \lambda_ {t} ^ {\prime}.\tag{3}
\]

This follows from Th. 2 of §4, No.~4, on taking into account the fact that \(t' \mapsto (t, t')\) is continuous and proper.

Remark. --- Lemmas 1,2,3 can be proved very simply without making use of the results of §§4 and 6, by a direct argument. For example, the relation (2) is obvious by definition if \(f \in \mathcal{K}(\mathrm{T} \times \mathrm{T}')\) . If \(f\) is lower semi-continuous on \(\mathrm{X} = \mathrm{T} \times \mathrm{T}'\) , it suffices to observe that \(t' \mapsto f_t(t')\) is the upper envelope of the functions \(t' \mapsto g_t(t') = g(t, t')\) , where \(g\) runs over the set of functions in \(\mathcal{K}(\mathrm{X})\) such that \(0 \leqslant g \leqslant f\) . Finally, for arbitrary \(f\) , one notes that if \(h \geqslant f\) is lower semi-continuous on \(\mathrm{X}\) , then \(t' \mapsto h(t, t')\) is lower semi-continuous on \(\mathrm{T}'\) ; and conversely, if \(t' \mapsto u(t')\) is lower semi-continuous on \(\mathrm{T}'\) and is such that \(u(t') \geqslant f(t, t')\) for all \(t' \in \mathrm{T}'\) , then the function \(h\) such that \(h(t, t') = u(t')\) , \(h(t_1, t') = +\infty\) for \(t_1 \neq t'\) , is lower semi-continuous on \(\mathrm{X}\) and satisfies \(h \geqslant f\) . Once Lemma 1 is proved, one deduces from it that the set \((\mathrm{T} - \{t\}) \times \mathrm{T}'\) is \(\lambda_t'\) -negligible, and it is then every easy to prove Lemmas 2 and 3.

The relation (3) permits denoting its two members by \(\int\mathbf{f}(t,t^{\prime})d\mu^{\prime}(t^{\prime})\) without risk of confusion. The analogous results obviously hold for the measures \(\lambda_{t^{\prime}}=\mu\otimes\varepsilon_{t^{\prime}}\) .

Instead of the notations

\[
\int^ {*} f (t, t ^ {\prime}) d \nu (t, t ^ {\prime}), \quad \int^ {\bullet} f (t, t ^ {\prime}) d \nu (t, t ^ {\prime}), \quad \int \mathbf {f} (t, t ^ {\prime}) d \nu (t, t ^ {\prime}),
\]

we shall employ the notations

\[
\iint^ {*} f (t, t ^ {\prime}) d \mu (t) d \mu^ {\prime} (t ^ {\prime}), \iint^ {\bullet} f (t, t ^ {\prime}) d \mu (t) d \mu^ {\prime} (t ^ {\prime}), \iint \mathbf {f} (t, t ^ {\prime}) d \mu (t) d \mu^ {\prime} (t ^ {\prime}),
\]

consistent with the notations adopted in Ch. III, §4, No.~1.

The interpretation of the measure \(\nu\) as an integral \(\int \lambda_t' d\mu(t)\) will allow us to translate the results of §3 into the language of product measures. On the other hand, the measure \(\lambda_t'\) is carried by \(\{t\} \times T' = \overline{\text{pr}}_1(t)\) , so that this integral defines a decomposition of \(\nu\) into slices, relative to the projection \(\text{pr}_1\) of \(T \times T'\) onto \(T\) (§6, No.~6). Before giving a list of the results so obtained, here is a useful property:

PROPOSITION 1. --- Let \((\mu_{\alpha})_{\alpha \in \mathrm{A}}\) (resp. \((\mu_{\beta}^{\prime})_{\beta \in \mathrm{B}})\) be a summable family of positive measures on T (resp. on \(\mathbf{T}'\) ), with sum denoted by \(\mu\) (resp. by \(\mu'\) ). The family \((\mu_{\alpha} \otimes \mu_{\beta}^{\prime})_{(\alpha, \beta) \in \mathrm{A} \times \mathrm{B}}\) is then summable on \(\mathbf{T} \times \mathbf{T}'\) , and

\[
\mu \times \mu^ {\prime} = \sum_ {(\alpha , \beta) \in \mathrm{A} \times \mathrm{B}} \mu_ {\alpha} \otimes \mu_ {\beta} ^ {\prime}.\tag{4}
\]

These properties are obvious when A and B are finite. It follows that if \(A'\) (resp. \(B'\) ) is a finite subset of A (resp. of B), then

\[
\sum_ {(\alpha , \beta) \in A ^ {\prime} \times B ^ {\prime}} \mu_ {\alpha} \otimes \mu_ {\beta} ^ {\prime} \leqslant \mu \otimes \mu^ {\prime}.
\]

The family \((\mu_{\alpha} \otimes \mu'_{\beta})\) is therefore summable. To show that the two members of (4) are equal, it suffices to prove that the second member satisfies the characteristic property of product measures (Ch. III, §4, No.~1, Th. 1), which is shown by the following calculation.

Let \(f\) be an element of \(\mathcal{H}_{+}(\mathrm{T})\) , \(f'\) an element of \(\mathcal{H}_{+}(\mathrm{T}')\) ; recall that \(f \otimes f'\) denotes the function \((t, t') \mapsto f(t)f'(t')\) on \(\mathrm{T} \times \mathrm{T}'\) , which belongs to \(\mathcal{H}_{+}(\mathrm{T} \times \mathrm{T}')\) (A, II, §7, No.~7). Then, by the definition of product measures,

\[
\begin{array}{r l} \sum_ {(\alpha , \beta) \in \mathrm{A} \times \mathrm{B}} \langle \mu_ {\alpha} \otimes \mu_ {\beta} ^ {\prime}, f \otimes f ^ {\prime} \rangle & = \sum_ {(\alpha , \beta) \in \mathrm{A} \times \mathrm{B}} \big (\langle \mu_ {\alpha}, f \rangle \langle \mu_ {\beta} ^ {\prime}, f ^ {\prime} \rangle \big) \\ & = \Big (\sum_ {\alpha \in \mathrm{A}} \langle \mu_ {\alpha}, f \rangle \Big) \Big (\sum_ {\beta \in \mathrm{B}} \langle \mu_ {\beta} ^ {\prime}, f ^ {\prime} \rangle \Big) \\ & = \langle \mu , f \rangle \langle \mu^ {\prime}, f ^ {\prime} \rangle \\ & = \langle \mu \otimes \mu^ {\prime}, f \otimes f ^ {\prime} \rangle . \end{array}
\]

